\documentclass[11pt]{article}
\usepackage[utf8]{inputenc}
\usepackage{amsmath, amssymb, amsthm}
\usepackage{geometry}
\geometry{margin=1in}

\newtheorem{theorem}{Theorem}
\newtheorem{lemma}[theorem]{Lemma}
\newtheorem{corollary}[theorem]{Corollary}
\theoremstyle{definition}
\newtheorem{definition}[theorem]{Definition}
\newtheorem{example}[theorem]{Example}

\title{Shannon Entropy}
\author{math-foundations / concept 36}
\date{}

\begin{document}
\maketitle

\subsection*{First, an example}
Before the general theory, fix intuition with a coin. Let $X \sim \mathrm{Bernoulli}(p)$,
so $X = 1$ with probability $p$ and $X = 0$ with probability $1-p$. The Shannon
entropy in bits is
\[
H(p) \;=\; -p \log_2 p - (1-p) \log_2 (1-p).
\]
\emph{Read aloud:} ``H of p equals minus p log-base-two of p, minus one-minus-p log-base-two of one-minus-p.''
At $p = 0.5$ the coin is fair and $H(0.5) = 1$ bit --- the maximum, because we are
maximally uncertain. At $p = 0.9$ the coin is biased and $H(0.9) \approx 0.469$ bits
--- less uncertain, since we can usually guess heads. At $p = 1$ the outcome is
deterministic and $H(1) = 0$ bits --- no surprise at all. Operationally, an optimal
Huffman code for i.i.d.\ draws of $X$ uses on average $H(p)$ bits per outcome, so
entropy is the price (in bits) of describing one sample.

\section{Motivation}
Self-information $I(x) = -\log p(x)$ measures the surprise of a single outcome.
Shannon entropy is the \emph{expected} self-information: the average surprise we feel
when drawing from $p$ repeatedly. Equivalently, it is the optimal expected length
(in bits, nats, hartleys, depending on the log base) of a prefix-free code for the
random variable.

\section{Definitions}

\begin{definition}[Shannon entropy, discrete]
Let $X$ be a discrete random variable on a countable alphabet $\mathcal{X}$ with
probability mass function $p$. The \emph{Shannon entropy} of $X$ is
\[
H(X) \;=\; \mathbb{E}\!\left[-\log p(X)\right] \;=\; -\sum_{x \in \mathcal{X}} p(x) \log p(x),
\]
\emph{Read aloud:} ``H of X equals the expectation of minus log p of X, equals minus the sum over x in script-X of p of x times log p of x.''
with the convention $0 \log 0 = 0$ (justified by
$\lim_{p \to 0^+} p \log p = 0$).
\end{definition}

\begin{definition}[Differential entropy, continuous]
Let $X$ be a continuous random variable with density $f$. The \emph{differential entropy}
of $X$ is
\[
h(X) \;=\; -\int_{\mathbb{R}} f(x) \log f(x) \, dx,
\]
\emph{Read aloud:} ``little-h of X equals minus the integral over the real line of f of x times log f of x, d x.''
whenever the integral exists. Unlike discrete entropy, $h(X)$ may be negative and is
not invariant under change of variables.
\end{definition}

The base of the logarithm fixes the units: $\log_2 \Rightarrow$ bits;
$\ln \Rightarrow$ nats; $\log_{10} \Rightarrow$ hartleys. We use $\log$ to denote a
generic base unless specified.

\section{Basic properties}

\begin{lemma}[Non-negativity]
For any discrete random variable $X$, $H(X) \ge 0$, with equality if and only if $X$ is
deterministic (i.e., some outcome has probability $1$).
\end{lemma}

\begin{proof}
Each summand $-p(x) \log p(x)$ is non-negative because $0 \le p(x) \le 1$ implies
$\log p(x) \le 0$. Equality in the sum requires every term to vanish, which forces
$p(x) \in \{0, 1\}$ for all $x$.
\end{proof}

\section{Maximum-entropy theorem}

The following result is the central inequality of this lesson.

\begin{theorem}[Uniform maximises entropy]
\label{thm:max-entropy}
Let $X$ be a discrete random variable on a finite alphabet $\mathcal{X}$ with
$|\mathcal{X}| = n$. Then
\[
H(X) \;\le\; \log n,
\]
\emph{Read aloud:} ``H of X is less than or equal to log n.''
with equality if and only if $X$ is uniform on $\mathcal{X}$ (i.e.,
$p(x) = 1/n$ for every $x \in \mathcal{X}$).
\end{theorem}

\begin{proof}
We use Jensen's inequality applied to the strictly concave function $\log$.
Let $\mathcal{X}^+ = \{x \in \mathcal{X} : p(x) > 0\}$ denote the support of $p$,
and write $m = |\mathcal{X}^+| \le n$. By definition,
\[
H(X) \;=\; -\sum_{x \in \mathcal{X}^+} p(x) \log p(x)
       \;=\; \sum_{x \in \mathcal{X}^+} p(x) \log \frac{1}{p(x)}.
\]
\emph{Read aloud:} ``H of X equals minus the sum over x in the support of p of x log p of x, equals the sum over x in the support of p of x times log of one over p of x.''
View the right-hand side as the expectation
$\mathbb{E}\!\left[\log \tfrac{1}{p(X)}\right]$, where the expectation is taken under $p$.
Since $\log$ is concave, Jensen's inequality gives
\[
\mathbb{E}\!\left[\log \tfrac{1}{p(X)}\right]
   \;\le\; \log\!\left(\mathbb{E}\!\left[\tfrac{1}{p(X)}\right]\right).
\]
\emph{Read aloud:} ``the expectation of log of one over p of X is less than or equal to log of the expectation of one over p of X.''
Compute the inner expectation:
\[
\mathbb{E}\!\left[\tfrac{1}{p(X)}\right]
   \;=\; \sum_{x \in \mathcal{X}^+} p(x) \cdot \tfrac{1}{p(x)}
   \;=\; \sum_{x \in \mathcal{X}^+} 1
   \;=\; m.
\]
\emph{Read aloud:} ``the expectation of one over p of X equals the sum over x in the support of p of x times one over p of x, equals the sum of ones over the support, equals m.''
Therefore
\[
H(X) \;\le\; \log m \;\le\; \log n,
\]
\emph{Read aloud:} ``H of X is less than or equal to log m, which is less than or equal to log n.''
where the second inequality uses $m \le n$ and monotonicity of $\log$.

For the equality condition: $\log$ is \emph{strictly} concave, so Jensen's inequality
is strict unless $1/p(X)$ is constant on the support, which means $p(x) = 1/m$ for all
$x \in \mathcal{X}^+$. Combined with $m \le n$, equality $H(X) = \log n$ requires $m = n$
and $p$ uniform on $\mathcal{X}$. Conversely, if $p$ is uniform on $\mathcal{X}$, then
$H(X) = -\sum_{x} (1/n) \log(1/n) = \log n$.
\end{proof}

\begin{corollary}
For any discrete random variable on $n$ outcomes, $0 \le H(X) \le \log n$. The lower
bound is achieved by deterministic $X$, and the upper bound by uniform $X$.
\end{corollary}

\section{Joint entropy and subadditivity}

\begin{definition}[Joint entropy]
For discrete random variables $X, Y$ with joint pmf $p(x, y)$,
\[
H(X, Y) \;=\; -\sum_{x, y} p(x, y) \log p(x, y).
\]
\emph{Read aloud:} ``H of X comma Y equals minus the double sum over x and y of p of x comma y times log p of x comma y.''
\end{definition}

\begin{theorem}[Subadditivity]
$H(X, Y) \le H(X) + H(Y)$, with equality iff $X$ and $Y$ are independent.
\end{theorem}

The proof uses non-negativity of KL divergence applied to the joint vs.\ product of
marginals; we defer it to the cross-entropy / KL lesson.

\section{Bernoulli entropy as a worked example}

\begin{example}[Binary entropy function]
For $X \sim \mathrm{Bernoulli}(p)$,
\[
H(p) \;=\; -p \log_2 p - (1-p) \log_2 (1-p).
\]
\emph{Read aloud:} ``H of p equals minus p log-base-two of p, minus one-minus-p log-base-two of one-minus-p.''
Differentiating, $H'(p) = \log_2\!\big(\tfrac{1-p}{p}\big)$, which vanishes at $p = 1/2$,
giving $H(1/2) = 1$ bit. By Theorem~\ref{thm:max-entropy} this is the global maximum, as
$\log_2 2 = 1$.
\end{example}

\end{document}
